\documentclass[a4paper]{article}

\hyphenpenalty 10000
\exhyphenpenalty 10000

\usepackage[T1]{fontenc}
\usepackage{enumitem}
\usepackage[
    top=0.79in,
    bottom=0.79in,
    left=0.79in,
    right=0.79in,
]{geometry}

\usepackage{xcolor}
\definecolor{blue}{HTML}{3A6EA5}
\usepackage{multicol}
\usepackage{titlesec}
\usepackage{enumitem}
\usepackage{hyperref}
\hypersetup{hidelinks}
\usepackage{array}

\pagestyle{empty}
\setlength{\parindent}{0pt}
\setlength{\tabcolsep}{10pt}
\newlength{\tablewidth}

\setlist[itemize]{topsep=0em, itemsep=0.0em}
\setlist[itemize,1]{before=\vspace{-1em}}

\newenvironment{period}[2]{%
\newcommand{\sarma}{#2}%
\setlength{\tablewidth}{\linewidth}
\addtolength{\tablewidth}{-2\tabcolsep}
\begin{tabular}{@{}p{0.85\tablewidth}>{\raggedleft\arraybackslash}p{0.15\tablewidth}@{}}%
#1 \newline
\begin{itemize}
}{%
\end{itemize} & \sarma \\%
\end{tabular}\\
}

\newenvironment{skills}{%
\setlength{\tablewidth}{\linewidth}
\addtolength{\tablewidth}{-2\tabcolsep}
\begin{tabular}{@{}p{0.15\tablewidth}p{0.85\tablewidth}@{}}
}{%
\end{tabular}
}
\titleformat{\section}{\color{teal}\uppercase}{}{0pt}{}[{\vspace{0.2em}\titlerule}]

\titleformat{\subsection}{\bf}{}{0pt}{}{}
\titlespacing{\section}{0pt}{0.25em}{0.25em}
\titlespacing{\subsection}{0pt}{0pt}{0.1em}

\begin{document}

\fontfamily{\sfdefault}
\selectfont

\begin{minipage}[t]{.5\textwidth}
{\LARGE \textbf{Marko Cvjetko}} \\[0.2em]
\textit{PhD Student}\\
Inria, FLOWERS team -- Bordeaux, France
\end{minipage}%
\begin{minipage}[t]{.5\textwidth}
\raggedleft
\href{mailto:marko.cvjetko@inria.fr}{\textcolor{blue}{marko.cvjetko@inria.fr}} \\
\href{https://github.com/markocvjetko}{\textcolor{blue}{github.com/markocvjetko}} \\
\href{https://markocvjetko.github.io/}{\textcolor{blue}{[personal website]}} \\
\href{https://scholar.google.co.uk/citations?hl=hr&user=WDyIW9QAAAAJ}{\textcolor{blue}{[Google Scholar]}}
\end{minipage}

\vspace{1em}

I am a 2nd year PhD student building curiosity-driven AI scientists that autonomously discover and control self-organizing phenomena, combining autotelic RL, vision--language models,   and evolutionary search. I am extending this program in two directions: integrating LLMs more deeply as discovery partners in the experimental loop, and testing these methods on more biologically grounded substrates through ongoing collaboration with the \href{https://drmichaellevin.org/}{\textcolor{blue}{Levin Lab}} (Tufts). I am a 2026 \href{https://diverseintelligencesummer.com/}{\textcolor{blue}{DISI}} fellow and will visit \href{https://www.epfl.ch/labs/csft/}{\textcolor{blue}{Cl\'{e}ment Hongler's group}} at EPFL in May 2026.
\vspace{0.4em}


\section{Publications}

\begin{description}[leftmargin=1.5em, style=unboxed, font=\normalfont\bfseries, itemsep=0.4em, parsep=0em, topsep=0.2em]
    \item[Under review] ~\\[0.2em]
    \hangindent=1em \textbf{The Artificial Experimentalist: Discovery and Control of Self-Organizing Phenomena with Autotelic Reinforcement Learning.} \underline{M. Cvjetko}, B. Hartl, M. Levin, C. Moulin-Frier, P.-Y. Oudeyer. Submitted to \textit{ALIFE 2026}. In collaboration with \href{https://drmichaellevin.org/}{\textcolor{blue}{Levin Lab, Tufts University}}. \href{https://developmentalsystems.org/carl/}{\textcolor{blue}{[project page]}}. \textit{Goal-conditioned RL agent that discovers and controls self-organizing patterns in cellular automata via minimal local interventions, with zero-shot generalization to unseen dynamics and action spaces.}
    \item[Journal (invited, in preparation)] ~\\[0.2em]
    \hangindent=1em Extension of an ALIFE 2025 paper for the \textit{Artificial Life} journal special issue ($\sim$5--8 papers selected from top-ranked submissions and best-paper nominees; 91 accepted at the conference).
    \item[Conference papers] ~\\[0.2em]
    \hangindent=1em \textbf{Exploring Flow-Lenia Universes with a Curiosity-driven AI Scientist: Discovering Diverse Ecosystem Dynamics.} T. Michel, \underline{M. Cvjetko} \textbf{(co-first authors)}, G. Hamon, P.-Y. Oudeyer, C. Moulin-Frier. In \textit{Proc. ALIFE 2025}, pp. 633--643. \href{https://doi.org/10.1162/ISAL.a.896}{\textcolor{blue}{doi:10.1162/ISAL.a.896}}. \href{https://developmentalsystems.org/Exploring-Flow-Lenia-Universes/}{\textcolor{blue}{[project page]}}. \textit{Population-based diversity search (IMGEPs) over system-wide complexity and evolutionary-activity metrics, surfacing ecosystem-level dynamics in Flow-Lenia.}\\[0.3em]
    \hangindent=1em \textbf{Expedition \& Expansion: Leveraging Semantic Representations for Goal-Directed Exploration in Continuous Cellular Automata.} S. Khajehabdollahi, G. Hamon, \underline{M. Cvjetko}, P.-Y. Oudeyer, C. Moulin-Frier, C. Colas. In \textit{Proc. ALIFE 2025}. \href{https://direct.mit.edu/isal/proceedings/isal2025/37/73/134068}{\textcolor{blue}{doi:10.1162/ISAL}}. \href{https://developmentalsystems.org/expedition_expansion/}{\textcolor{blue}{[project page]}}. \textit{Hybrid exploration alternating novelty search with VLM-generated linguistic goals, providing evidence that LLM-generated goals act as stepping stones into unexplored regions of the behavioral space.}
    \item[Workshop paper] ~\\[0.2em]
    \hangindent=1em \textbf{Evolving Symbiosis, from Barricelli's Legacy to Collective Intelligence: a Simulated and Conceptual Approach.} J. Ashford, \underline{M. Cvjetko}, R. L\"{o}ffler, B. Sakallioglu, A. Valerio, M. Tataryn, B. Hartl, L. Pio-Lopez, S. Nichele. Accepted to \textit{GECCO 2026 Workshop Proceedings}. \href{https://arxiv.org/abs/2603.08463}{\textcolor{blue}{arXiv:2603.08463}}. Runner-up award at \href{https://aliceworkshop.org/}{\textcolor{blue}{ALICE 2026 Workshop}}, Copenhagen. \textit{Revisits Barricelli's symbiogenesis simulations as a substrate for studying the emergence of cooperation and collective intelligence.}
    \item[Late-breaking abstract] ~\\[0.2em]
    \hangindent=1em \textbf{Discovering and Controlling Diverse Self-Organised Patterns in Cellular Automata Using Autotelic Reinforcement Learning.} \underline{M. Cvjetko}, G. Hamon, C. Moulin-Frier, P.-Y. Oudeyer. \textit{ALIFE 2025}, Kyoto, Japan. \href{https://hal.science/hal-05399881v1/document}{\textcolor{blue}{[HAL]}}. \href{https://developmentalsystems.org/rl-for-ca/}{\textcolor{blue}{[project page]}}
\end{description}


\section{Honors, Fellowships \& Research Visits}

\begin{period}{\textbf{\href{https://diverseintelligencesummer.com/}{\textcolor{blue}{Diverse Intelligences Summer Institute (DISI) Fellow}}}}{June 2026}
    \item Selected for a 3-week transdisciplinary residency on the origins, nature, and future of intelligences (cognitive science, AI, biology, philosophy); cohort of $\sim$40 fellows funded by the John Templeton Foundation.
\end{period}

\begin{period}{\textbf{Funded research visit, \href{https://www.epfl.ch/labs/csft/}{\textcolor{blue}{group of Cl\'{e}ment Hongler (EPFL)}}}}{May 2026}
    \item Two-week visit funded by the host group, coinciding with the \href{https://www.dem.eco/}{\textcolor{blue}{Demeco 2026}} workshop.
\end{period}

\begin{period}{\textbf{\href{https://aliceworkshop.org/}{\textcolor{blue}{ALICE Workshop}} Runner-up Award, Copenhagen}}{Feb 2026}
    \item Awarded for ``Evolving Symbiosis, from Barricelli's Legacy to Collective Intelligence''.
\end{period}


\section{education}

\subsection{\textbf{University of Bordeaux, Inria, FLOWERS team, France}}
\begin{period}{\textbf{PhD thesis -- Studying the emergence of open-ended evolution in cellular automata using curiosity-driven AI}}{Oct 2024 -- current}
    \item  \textbf{Supervised by \href{http://www.pyoudeyer.com/}{\textcolor{blue}{Pierre-Yves Oudeyer}} and \href{}{\textcolor{blue}{Cl\'{e}ment Moulin-Frier}}}
    \item Applying curiosity-driven AI (evolutionary search, autotelic RL, LLM-driven exploration) to the automated discovery and control of emergent phenomena in Artificial Life systems such as \href{https://www.youtube.com/watch?v=HT49wpyux-k}{\textcolor{blue}{Lenia}}.
    \item \textbf{Supervision}: co-supervising a Master's intern on the emergence of memory and learning in cellular automata, and two BSc students on a project investigating Turing-completeness in Lenia.
    \item \textbf{Service}: reviewer for \textit{ALIFE 2026}.
    \item \textbf{Collaborations}: ongoing collaboration with Levin Lab (Tufts University) on autotelic agents for the discovery and control of self-organizing patterns; upcoming 2-week funded research visit hosted by \href{https://www.epfl.ch/labs/csft/}{\textcolor{blue}{Cl\'{e}ment Hongler (EPFL)}}, coinciding with the \href{https://www.dem.eco/}{\textcolor{blue}{Demeco 2026}} workshop (May 2026).
\end{period}


\subsection{\textbf{University of Zagreb, Faculty of Electrical Engineering and Computing (FER)}, Croatia}
\begin{period}{\textbf{\href{https://www.fer.unizg.hr/en/studies/cro/master/fer3/computing/cs#}{\textcolor{blue}{MSc, Computer Science}}}}{Sep 2021 -- Jul 2024}
    \item \textbf{MSc thesis}: Recognition of the respiratory signal from thermal video.
        \begin{itemize}
                \item[] \textbf{Supervised by \href{https://haeslerlab.sites.vib.be/en#/}{\textcolor{blue}{Prof. Sebastian Haesler}} and \href{https://www.zemris.fer.hr/~ssegvic/index_en.html}{\textcolor{blue}{Prof. Sini\v{s}a \v{S}egvi\'{c}}}}
              \item[] In collaboration with \href{https://nerf.be/en#/}{\textcolor{blue}{Neuro-Electronics Research Flanders}}, project described in \textit{\textcolor{teal}{work experience}} section.
          \end{itemize}
    \item \textbf{Selected courses}:
        \textit{\href{https://www.fer.unizg.hr/en/course/maclea1}{\textcolor{blue}{Machine Learning}}},
        \textit{\href{https://www.fer.unizg.hr/en/course/deelea1_a}{\textcolor{blue}{Deep Learning}}},
        \textit{\href{https://www.fer.unizg.hr/en/course/sofcom}{\textcolor{blue}{Soft Computing}}},
        \textit{\href{https://www.fer.unizg.hr/en/course/parpro_b}{\textcolor{blue}{Parallel Programming}}}

\end{period}
\vspace{-0.1em}
\begin{period}{\hspace{12} \textbf{Erasmus student exchange at KU Leuven}, Flanders, Belgium}{Sep 2022 -- Jul 2023}
\item Coursework from the \href{https://www.kuleuven.be/programmes/master-artificial-intelligence}{\textcolor{blue}{Advanced Master AI programme}}
\item \textbf{Selected courses}:
        \textit{\href{https://onderwijsaanbod.kuleuven.be/syllabi/e/H02B2AE.htm#activetab=doelstellingen_idp1064224}{\textcolor{blue}{Cognitive Science}}},
        \textit{\href{https://onderwijsaanbod.kuleuven.be/syllabi/e/W0W26AE.htm#activetab=doelstellingen_idm10328960}{\textcolor{blue}{Philosophy of Mind}}},
        \textit{\href{https://onderwijsaanbod.kuleuven.be/2023/syllabi/e/H02B6AE.htm#activetab=doelstellingen_idp955552}{\textcolor{blue}{Linguistics and AI}}}
        \textit{\href{https://onderwijsaanbod.kuleuven.be/syllabi/e/H02D2AE.htm#activetab=doelstellingen_idm323056}{\textcolor{blue}{Uncertainty in AI}}},
        \textit{\href{https://onderwijsaanbod.kuleuven.be/2023/syllabi/e/H02B3AE.htm#activetab=doelstellingen_idp1044304}{\textcolor{blue}{Neural Computing}}},

\end{period}
\vspace{-0.15em}
\begin{period}{\textbf{\href{https://www.fer.unizg.hr/en/studies/bachelor/computing}{\textcolor{blue}{BSc, Computing}}}}{Sep 2016 -- Sep 2021}
    \item \textbf{BSc thesis}: \textit{Music Synthesis using Artificial Intelligence} -- LSTM next-note prediction on symbolically represented music. Supervised by \href{https://www.zemris.fer.hr/~yeti/}{\textcolor{blue}{Prof. Domagoj Jakobovi\'{c}}}.
\end{period}


\section{Work Experience}

    \begin{period}{\textbf{Graduate Researcher, \href{https://nerf.be/en#/}{\textcolor{blue}{Neuro-Electronics Research Flanders}} (\href{https://www.imec-int.com/en}{\textcolor{blue}{imec}}/KU Leuven/VIB), \href{https://haeslerlab.sites.vib.be/en#/}{\textcolor{blue}{Haesler Lab}}} (\href{https://github.com/}{\textcolor{blue}{GitHub}})}{Mar 2023 -- Jul 2024}
        \item \textbf{Supervised by \href{https://haeslerlab.sites.vib.be/en#/}{\textcolor{blue}{Prof. Sebastian Haesler}}}
        \item Developed a machine learning framework for extracting and analyzing respiratory responses of head-restrained mice from infrared footage, in collaboration with experimental neuroscientists at Haesler Lab.
        \item Tackled low-data and domain-shift challenges via unsupervised and self-supervised learning with vision transformers and deep CNNs.
        \item Configured the data annotation platform, oversaw annotation tasks and trained students on annotation guidelines. Presented in a \href{https://www.figma.com/design/SPusXHMO6yUeBIAR5LSxHE/retreat(1)?node-id=0%3A1&t=jTCQF1uwkTbBV0Az-1}{\textcolor{blue}{poster session}} at the institute's yearly retreat.
    \end{period}

    \begin{period}{\textbf{Machine Learning Engineer Internship, \href{https://site.doxray.com/}{\textcolor{blue}{DoXray}}}}{Jun 2022 -- Sep 2022}
            \item Fine-tuned LayoutXLM --- a multilingual multimodal (text + layout + vision) transformer --- for named-entity recognition on German invoice documents, benchmarked against a BERT baseline. \href{https://blog.doxray.com/p/named-entity-recognition-using-a}{\textcolor{blue}{Blogpost}}.
    \end{period}


\section{Projects}


    \begin{period}{\textbf{Generative models for unsupervised behavioral spaces in open-ended discovery}}{2026 -- current}
        \item Training diffusion models and variational autoencoders to represent complex system behavior and sample their parameters for open-ended diversity search
        \item Goal: replace hand-designed semantic descriptors with learned, generative behavior spaces that enable scalable goal-directed search and quality--diversity over continuous cellular automata.
    \end{period}

    \begin{period}{\textbf{Music-reactive cellular automata}}{Nov 2024}
        \item Real-time music-driven Lenia animations: tempo and frequency features drive a continuous cellular automaton \href{https://www.youtube.com/watch?v=M5BxJMfp3uk}{\textcolor{blue}{(example)}}. Built during the \href{https://sites.google.com/view/hack1robo/accueil}{\textcolor{blue}{hack1robo}} hackathon.
    \end{period}

    \begin{period}{\textbf{Cellular Automata Study in Jax} (\href{https://github.com/}{\textcolor{blue}{GitHub}})}{Jan -- Mar 2024}
        \item Lenia implementation in Jax; trained a 13-parameter neural network to learn Conway's Game of Life update rule.
    \end{period}

    \begin{period}{\textbf{Automatic Universal Taxonomies for Multi-Domain Semantic Segmentation} (\href{https://github.com/}{\textcolor{blue}{GitHub}})}{2023 -- 2024}
        \item Joint label space for semantic segmentation datasets via cross-dataset confusion matrices, enabling to train a single classification head on the union.
    \end{period}

    \begin{period}{\textbf{Can Augmentation Make Better Doctors? Data Augmentation in a Low-Resource Sequence Labeling Task}}{2022}
        Benchmarked data augmentations for BERT-based sequence labeling on a low-resource medical named entity recognition dataset.
    \end{period}


\vspace{-0.23em}
\section{skills}
\begin{skills}
    \textbf{Programming} & Python (Jax, PyTorch, NumPy, Transformers, Optuna, \ldots), C, Bash \\
    \\
    \textbf{Tooling} & Linux, Slurm, Docker, Singularity, Git, LaTeX \\
    \\
    \textbf{Languages} & Croatian (native), English (fluent), French (intermediate) \\
\end{skills}

\end{document}
